
%\newpage
\section{Selbsttest Funktion}
\label{sec:selftest}
Beginnend mit Version 0.9k habe ich eine Selbsttest-Funktion eingebaut.
Die Benutzung ist sehr einfach.
Wenn Sie Pr\"ufspitzen mit Klemmen verwenden, klemmen Sie alle Klemmen auf ein unisoliertes St\"uck Draht und dr\"ucken Sie
den Start Taster.
Das Programm bemerkt die kurzgeschlossenen Klemmen und startet den Selbsttest.
Nach Beenden des Selbsttestes wird der Tester mit normaler Messung fortfahren.
Wenn kein Bauteil angeschlossen ist, wird der TransistorTester mit der
Meldung ,,kein,unbek. oder defektes Bauteil'' enden.
Die ungl\"uckliche Seite der Selbsttest-Funktion ist, da"s der 8k flash Speicher fast verbraucht ist.
Die L\"ange der ATmega8 Version ist ungef\"ahr 8000 Bytes.
Man kann die Selbsttest-Funktion zusammen mit allen anderen Optionen nur f\"ur einen ATmega168
oder ATmega328 w\"ahlen.
F\"ur einen ATmega8 mu"s man wegen des begrenzten flash Speichers wenigstens die Option AUTOSCALE\_ADC weglassen.
Die einzelnen Schritte der Selbsttestfunktion werden immer in Zeile~1 des LCD Displays mit dem Buchstaben T
gefolgt von der Schritt-Nummer dargestellt.
Jeder Schritt wird achtmal wiederholt, bevor das Programm zum n\"achsten Schritt geht.
Aber wenn man den Start Taster gedr\"uckt h\"alt, wenn ein Durchlauf beendet wird, wird dieser Test nicht
mehr wiederholt.
Wenn man den Taster w\"ahrend dem gesamten Selbsttest gedr\"uckt h\"alt, wird jeder Schritt nur einmal ausgef\"uhrt.

In jedem Schritt werden nur Me"sergebnisse dargestellt, es wird keine Fehleranalyse durchgef\"uhrt.
Man  mu"s die Ergebnisse selbst interpretieren.
An dieser Stelle m\"ochte ich noch einen wichtigen Hinweis geben. F\"uhren Sie niemals Messungen mit
eingestecktem ISP-Stecker durch!
Die ISP-Schnittstelle st"ort die Messungen.

\vspace{1cm}
Hier ist die Liste der derzeit eingebauten Tests:
\vspace{1cm}
\begin{enumerate}
\item {\bf Messung der 1,3V (oder 1,1V) Referenz Spannung (band gap Reference).} In Zeile~1 wird der Text ,,Ref='' und die
gemessene Spannung in mV angezeigt.
F\"ur den ATmega8 sollte die gemessene Spannung nahe bei 1,3V liegen, f\"ur die anderen Prozessoren liegt
die Referenzspannung normal um die 1,1V.
Die zweite Zeile zeigt den resultierenden Faktor f\"ur die Kapazit\"ats-Messung mit dem \(470k\Omega\) Widerstand.
\item {\bf Vergleich der \(680\Omega\) Widerst\"ande.} 
In Zeile~1 wird der kryptische Text  ,,+RL- 12 13 23'' angezeigt. Das bedeutet das folgende:
Das RL ist eine Kurzform f\"ur Resistor Low, was die \(680\Omega\) Widerst\"ande meint. Die ,,12'' steht f\"ur: 
Widerstand am Pin~1 ist verbunden mit VCC (+) und Widerstand am Pin~2 ist verbunden mit GND (-). 
Das Ergebnis dieser Messung steht in Zeile~2 an erster Stelle.
 In Zeile~1 folgt nun ,,13'', was bedeutet, da"s der Widerstand von Pin~1 wieder nach VCC verbunden ist,
aber da"s jetzt der \(680\Omega\) Widerstand von Pin~3  mit GND verbunden ist.
Das Ergebnis steht in Zeile 2 an der mittleren Stelle.

Die letzte Messung dieses Tests ,,23'' bedeutet, da"s nun der Widerstand von Pin~2 mit VCC verbunden ist und
der Widerstand von Pin~3 mit GND verbunden ist.
Das Ergebnis steht in der zweiten LCD-Zeile an letzter Stelle.
Ich m\"ochte daran erinnern, da"s die Aufl\"osung des ADC etwa 4,88mV betr\"agt!
Die Me"ssituation wird auch in Abbildung~\ref{fig:test2} dargestellt.
Alle diese Kombinationen sollten im Hinblick auf die internen Port-Widerst\"ande folgendes ergeben:
\(\frac{5001 \cdot  (19+680)}{ (19+680+680+22)} = 2493\) .

\begin{figure}[H]
\centering
\includegraphics[width=17cm]{../FIG/Test2.eps}
\caption{Vergleich der \(680\Omega\) Widerst\"ande }
\label{fig:test2}
\end{figure}

\item {\bf Vergleich der \(470k\Omega\) Widerst\"ande.}
Jetzt zeigt das Display in Zeile 1 ,,+RH- 12 13 23''. Das gleiche Verfahren wie in Schritt wird wiederholt mit den
 \(470k\Omega\) Widerst\"anden (Symbol RH).
Die Ergebnisse sollten nahe bei \(\frac{5001 \cdot (19 + 470000]}{ (19 + 470000 + 470000 + 22)} = 2500\) f\"ur alle Kombinationen sein.
\item In diesem Schritt wird nichts gemessen, es erfolgt lediglich die {\bf Anweisung ,,isolate Probe''},
was bedeutet, da"s es nun Zeit ist die Klemmen zu trennen (vom blanken Draht l\"osen).
\item Dieser Schritt testet die {\bf F\"ahigkeit von mit GND (-) verbundenen \(470k\Omega\) Widerst\"anden (H) die Testpins auf GND zu ziehen.}
Zeile 1 zeigt den Text ,,RH-''.
Zeile 2 sollte f\"ur alle drei Pinne  Null mV anzeigen.
\item Dieser Schritt testet die {\bf F\"ahigkeit von mit VCC (+) verbundenen \(470k\Omega\) Widerst\"anden (H) die Testpins auf VCC zu ziehen.}
Zeile 1 zeigt den Text ,,RH+''.
Der bestm\"ogliche mV Wert f\"ur die drei Messungen w\"are 5001 in Zeile~2.
 Gro"se Abweichungen vom Idealwert f\"ur Schritt 5 und 6 sind Fehler wie Isolations Problem, Flu"smittelreste oder besch\"adigter Port.
\item {\bf Messung des Innenwiderstandes eines auf GND geschalteten Ausgangs.}
Der Text in der ersten Zeile des LCD ist ,,Ri\_Lo = (mV)''.
In der zweiten Zeile des LCDs werden drei Spannungswerte angezeigt.
Die internen Port C Widerst\"ande von auf GND (-) geschalteten Ausg\"angen werden mit dem Strom
der nach VCC (+) geschalteten \(680\Omega\) Widerst\"ande gemessen, siehe Abbildung~\ref{fig:test7}.
Es werden nur die drei Pinne des ADC Ports vermessen, die Widerstands-Ports  PB0, PB2 und PB4 k\"onnen nicht
ohne Ver\"anderung der Hardware gemessen werden.
Es wird angenommen, da"s die Port Widerst\"ande von unterschiedlichen Ports nahezu identisch sind.
Um die Widerstandswerte zu erhalten, m\"ussen die angezeigten mV Werte durch ungef\"ahr 7 geteilt werden (siehe Schritt 8).
\begin{figure}[H]
\centering
\includegraphics[]{../FIG/Test7.eps}
\caption{Messung des Innenwiderstandes von auf GND geschalteten Port C Ausg\"angen.}
\label{fig:test7}
\end{figure}

\item {\bf Messung des Innenwiderstandes von auf VCC geschalteten Port Ausg\"angen.}
Der ben\"otigte Strom wird von auf GND geschalteten \(680\Omega\) Widerst\"anden geliefert.
Der Text in der ersten LCD Zeile ist ,,Ri\_Hi= (mV)'', wenn keine AUTO\_CAL Option gesetzt ist.
Mit gesetzter AUTO\_CAL Option wird nur ,,Ri\_Hi='' ausgegeben.
In der zweiten Zeile des LCDs werden drei Spannungswerte angezeigt als Differenz zu VCC.
Es ist die gleiche Messung wie die in Test~7 zur anderen Seite, wie in Abbildung \ref{fig:test8} gezeigt wird.
Mit folgenden Schritten k\"onnen Sie den Innenwiderstand selbst berechnen:
Um den Strom zu berechnen: \((5001 - (Ergebnis von Test 7) - (Ergebnis von Test 8)) / 680\).
Sie k\"onnen beide Widerstandswerte erhalten, indem Sie die jeweilige Spannung
(Ergebnisse von Test~7 oder Test~8) durch diesen Strom dividieren.
Wenn die AUTO\_CAL Option gew\"ahlt wurde und der vorige und dieser Test achtmal wiederholt wurden,
wird der Mittelwert der Innenwiderst\"ande der Ausgabe-Ports im EEprom festgehalten.
Das Ergebnis f\"ur diesen Test wird dann in Zeile~1 in \(d\Omega\) angezeigt, der letzte Wert der Zeile~2 ist dann
das Ergebnis f\"ur den vorigen Test (auf GND geschaltet).
Diese Werte stehen dann f\"ur weitere Messungen zur Verf\"ugung.
\begin{figure}[H]
\centering
\includegraphics[]{../FIG/Test8.eps}
\caption{Messung des Innenwiderstandes des auf VCC geschalteten Port C Ausg\"angen.}
\label{fig:test8}
\end{figure}

\item {\bf Messung des Nulloffsets der Kondensatormessung.}
F\"ur die Pinkombinationen 1:3, 2:3 und 1:2 wird der Nullwert der Kondensatormessung in pF ausgegeben.
In der Software wird ein Vorgabewert von circa 39 pF f\"ur die normale Me"swertausgabe ber\"ucksichtigt.
F\"ur die Ausgabe dieses Tests ist keine Korrektur ber\"ucksichtigt, es wird auch kein Nulloffset abgezogen.
Wenn die AUTO\_CAL Option gew\"ahlt wurde und der Test achtmal wiederholt wurde, wird der Mittelwert
der gefundenen Nulloffsets im EEprom festgehalten und steht f\"ur die weiteren Kapazit\"ats-Messungen zur Verf\"ugung.
Angezeigt wird das durch die Ausgabe ,,EE=xx''  in Zeile 1.
Bitte beachten Sie, da"s Ver\"anderungen des Me"saufbaus einen Neuabgleich sinnvoll machen.
So kann sich der Nulloffset durch Verwenden von Kabeln mit Klemmen um circa 3~pF gegen\"uber einem leeren
Sockel erh\"ohen.

\item {\bf Messung des Komparatoroffsets f\"ur den Abgleich der Kondensatormessung.}
Um den Offset des analogen Komparators zu bestimmen, mu"s ein Kondensator an Pin~1 und Pin~3 angeschlossen werden.
Der Kondensator wird f\"ur die Pufferung der Ladespannung bei der Kondensatormessung ben\"otigt, um den Unterschied der 
Ladespannung zur internen Referenzspannung bestimmen zu k\"onnen. Es sollte ein Kondensator hoher G\"ute mit
mindestens \(100 nF\) sein. Kapazit\"atswerte von bis zu \(22 \mu F\) k\"onnen problemlos verwendet werden.
Auf keinen Fall sollte man Elektrolyt-Kondensatoren verwenden.
Bei erfolgreicher Messung wird der Korrekturwert kurz in Zeile 2 angezeigt und in den EEprom Speicher geschrieben.

\end{enumerate}

An Ende der Selbsttest-Funktion wird der Text ,,Auto Test End'' in Zeile~1 und die Versions-Nummer der Software in Zeile~2 angezeigt.
Wenn gew\"ahlt, wird ein {\bf 50 Hz Rechtecksignal} auf Testpin~2 generiert und das gegenphasige Signal auf Testpin~3.
Testpin~1 wird auf GND geschaltet. Der Strom f\"ur Testpin 2 und 3 wird mit \(680\Omega\) Widerst\"anden begrenzt.
Angezeigt wird das durch die Anzeige 50Hz am Ende der Zeile 1.
Diese Ausgabe wird dreizigmal mit jeweils 2~Sekunden Dauer durchgef\"uhrt.
Man kann die Zeit der Verz\"ogerungs-Aufrufe pr\"ufen, wenn man ein Oszilloskop oder einen
Frequenz-Z\"ahler besitzt.
Wenn man f\"ur die Takterzeugung keinen Quarz benutzt, k\"onnen die
Ergebnisse der Kondensatormessung ungenau sein.
Eine genaue Taktfrequenz und genaue Wartezeiten sind wichtig f\"ur die Bestimmung von Kapazit\"atswerten.
Die Ausgabe des 50Hz Signals kann durch gedr\"uckthalten des Start-Tasters vorzeitig abgebrochen werden.

Dann f\"ahrt das Programm mit der normalen Me"sfunktion fort.

\begin{table}[H]
  \begin{center}
    \begin{tabular}{| l | c | c | c |}
    \hline
Test Nr. & 1. Ergebnis & 2. Ergebnis & 3. Ergebnis \\
    \hline
    \hline
Test 1 & band gap Ref  1228 &  & \\
Referenz  & RHfakt 756 &  &  \\
    \hline
Test 2 & RL1+ RL2- & RL1+ RL3- & RL2+ RL3- \\
Vergleich \(680\Omega\) & 2489 & 2488 & 2488 \\
bester Wert: 2493 & & & \\
    \hline
Test 3 & RH1+ RH2- & RH1+ RH3- & RH2+ RH3- \\
Vergleich \(470k\Omega\) & 2493 & 2493 & 2493 \\
bester Wert: 2500 & & & \\
    \hline
Test 4 & isolate probe & & \\
    \hline
Test 5 & RH1- &  RH2- & RH3- \\
\(470k\Omega\) isolation & 0 & 0 & 0 \\
bester Wert: 0 & & & \\
    \hline
Test 6 & RH1+ & RH2+ & RH3+ \\
\(470k\Omega\) isolation & 4996 & 4995 & 4995 \\
bester Wert: 5001 & & & \\
    \hline
Test 7 & TP1- RL1+ & TP2- RL2+ & TP3- RL3+ \\
Pin Widerstand Low & 132 & 132 & 137 \\
    \hline
Test 8 & TP1+ RL1- & TP2+ RL2- & TP3+ RL3- \\
Pin Widerstand High & 151 & 151 & 151 \\
    \hline
Test 9 & C0  1:3 & C0  2:3 & C0 1:2 \\
Kondensator Offset & 34 & 35 & 33 \\
    \hline
    \end{tabular}
  \end{center}
  \caption{Selbsttest Mega8, Signature 1E9307, WITH\_AUTOREF}
  \label{tab:test_m8} 
\end{table}

\begin{table}[H]
  \begin{center}
    \begin{tabular}{| l | c | c | c |}
    \hline
Test No. & 1. Ergebnis & 2. Ergebnis & 3. Ergebnis \\
    \hline
    \hline
Test 1 & band gap Ref  1051 &  & \\
Referenz  & RHfakt 902 &  &  \\
    \hline
Test 2 & RL1+ RL2- & RL1+ RL3- & RL2+ RL3- \\
Vergleich \(680\Omega\) & 2493 & 2493 & 2493 \\
bester Wert: 2493 & & & \\
    \hline
Test 3 & RH1+ RH2- & RH1+ RH3- & RH2+ RH3- \\
Vergleich \(470k\Omega\) & 2494 & 2498 & 2498 \\
bester Wert: 2500 & & & \\
    \hline
Test 4 & isolate probe & & \\
    \hline
Test 5 & RH1- &  RH2- & RH3- \\
\(470k\Omega\) isolation & 1 & 1 & 1 \\
bester Wert: 0 & & & \\
    \hline
Test 6 & RH1+ & RH2+ & RH3+ \\
\(470k\Omega\) isolation & 5000 & 5000 & 5000 \\
bester Wert: 5001 & & & \\
    \hline
Test 7 & TP1- RL1+ & TP2- RL2+ & TP3- RL3+ \\
Pin Widerstand Low & 138 & 139 & 139 \\
    \hline
Test 8 & TP1+ RL1- & TP2+ RL2- & TP3+ RL3- \\
Pin Widerstand High & 156 & 156 & 156 \\
    \hline
Test 9 & C0  1:3 & C0  2:3 & C0 1:2 \\
Kondensator Offset & 39 & 39 & 38 \\
    \hline
    \end{tabular}
  \end{center}
  \caption{Selbsttest Mega168, Signature 1E9406, WITH\_AUTOREF}
  \label{tab:test_m168} 
\end{table}
